% Generated draft: requires mathematical and source review.
\documentclass[11pt]{article}
\usepackage[margin=25mm]{geometry}
\usepackage{fontspec}
\setmainfont{Latin Modern Roman}
\usepackage{amsmath,amssymb,mathtools,unicode-math}
\setmathfont{Latin Modern Math}
\usepackage{enumitem,needspace,xurl,hyperref}
\setlength{\parindent}{0pt}
\setlength{\parskip}{5pt}
\setlength{\emergencystretch}{3em}
\newcommand{\N}{\mathbb{N}}
\newcommand{\Z}{\mathbb{Z}}
\newcommand{\Q}{\mathbb{Q}}
\newcommand{\R}{\mathbb{R}}
\newcommand{\C}{\mathbb{C}}
\newcommand{\F}{\mathbb{F}}
\newcommand{\A}{\mathbb{A}}
\newcommand{\PP}{\mathbb{P}}
\newcommand{\E}{\mathbb{E}}
\newcommand{\eps}{\varepsilon}
\newcommand{\dd}{\,\mathrm{d}}
\begin{document}
\section*{UnsolvedMath ID: 30006482}
\textbf{Source problem code:} catalogue entry
\subsection{Definitions and mathematical statement}
The supplied record, entitled ``Huang--Yang Completion-of-Squares Fragment,'' does not determine a standalone mathematical conjecture. Its supplied statement status is \texttt{unrecoverable}. The supplied verification note identifies the extraction as a technical proof fragment, rather than an open problem; this qualification prevents treating its leading word ``conjecture'' as a recoverable conjecture statement.

The one explicit operator convention in the fragment is
\[
|A|^{2}:=A^{*}A.
\]
For an operator on a Hilbert space, this gives
\[
\langle\psi,A^{*}A\psi\rangle=\|A\psi\|^{2}\geq 0
\]
whenever the expressions are defined. For unbounded operators, the corresponding quadratic form is defined on the domain of \(A\). The source does not specify the Hilbert space or operator domains relevant to its Hamiltonian.

The fragment describes a completion-of-squares argument involving terms denoted by \(H_0,Q_2,Q_3,Q_4\) and corrections denoted by \(T_{\sigma}(k)\) and \(T_{\sigma,\sigma'}(x,y)\). It reports that a standard Cauchy--Schwarz estimate bounds the expectation of \(Q_3\) by a quantity of order \(O(L^{3}\rho^{7/3})\), whereas its expectation on an unspecified trial state is zero. It therefore proposes introducing the corrections, decomposing the Hamiltonian into squares, and dropping nonnegative terms to obtain a lower bound.

The extracted display numbered \((3)\) begins with \(H_0+Q_2+Q_4\) and an approximation symbol. Its readable components include square terms, a term involving \(\rho_{\uparrow}\rho_{\downarrow}L^{3}\), \(V\), and \(\varphi^{2}\), the subtraction \(-L^{3}C_L(\rho_{\uparrow},\rho_{\downarrow})\), a spatial integral with weight \(V(x-y)\), a sum restricted by \(\sigma\ne\sigma'\), and a remainder written \(O(\rho N)\). The extraction does not reliably preserve their complete arrangement, kinetic coefficient, or operator notation, so it cannot be promoted to an identity or inequality.

No precise quantified mathematical claim can be recovered from these data. In particular, the corrections, Hamiltonian terms, state class, spatial and momentum domains, boundary conditions, potential assumptions, asymptotic regime, meaning of the approximation, and target lower bound are unspecified. Nonnegativity of an integral weighted by \(V\) would also require hypotheses absent from the record. These unresolved qualifications must be supplied from a reliable complete formulation before a statement can be formalized.
\subsection{Short English statement}
This record preserves a damaged fragment of a Huang--Yang completion-of-squares argument. It describes using nonnegative operator squares to obtain a Hamiltonian lower bound, but supplies neither the hypotheses nor a recoverable complete formula. Although the registry labels it open, its supplied verification note identifies it as a proof fragment rather than an independent open problem. A precise conjecture cannot be stated from the supplied evidence.
\subsection{Sources}
\begin{itemize}
\item Local UnsolvedMath registry, supplied complete source record for permanent ID \(30006482\), source problem code \texttt{catalogue entry}, title ``Huang--Yang Completion-of-Squares Fragment.'' Registry paths supplied by the project are \texttt{data/unsolvedmath/problems.json} and \texttt{catalog/problems.json}; their contents have not been independently inspected here. The record belongs to the Oberwolfach Reports Open Problems collection and supplies no named proposer.
\item Supplied source citation: \emph{Arbeitsgemeinschaft: Analysis of Many-body Quantum Systems} (2026), Oberwolfach Reports, pp. 2491--2544, DOI \url{https://doi.org/10.4171/owr/2025/46}. The extracted page marker is 2529. The supplied literature-source label is ``Oberwolfach Report 46/2025.'' The difference between the citation year 2026 and the report identifier and proposed year 2025 is unresolved.
\item Supplied additional literature source: \emph{Rigorous upper bound for the Gell--Mann--Brueckner formula}, \url{https://doi.org/10.1007/s00220-023-04672-2}. Its precise relevance to the extracted formula is not established by the supplied record.
\item Literature assessment dated exactly \texttt{2026-08-22}: the supplied curation labels the extraction open because it regards it as an unresolved intermediate completion-of-squares estimate in the Huang--Yang analysis. This is supplied dated curation, not an independently verified assessment. The background explicitly requests independent verification of time-sensitive status and cited claims.
\item The supplied statement-verification note says that the official report identifies the extraction as a technical proof fragment, not an open problem, and that no separate clean statement is promoted. This report-verification claim is attributed to the supplied record and has not been independently checked here.
\end{itemize}
\end{document}
